\begingroup
\setstretch{1.1}
\section{Supplementary evaluation}
\label{app:evaluation}

The test-set counts in Table~\ref{tab:heldout_counts} use the same ensemble as the main catalogue. Each reference class contains three test sources; the source-level predictions and full confusion matrix are retained with the accompanying analysis outputs.

\begin{table}[H]
\centering
\small
\setstretch{1}
\begin{tabular}{@{}lrrl@{}}
\toprule
Reference class & Test sources & Correct & Other assignments \\
\midrule
Amorphous & 3 & 3 & -- \\
ACJ & 3 & 1 & Spirals (1), Structured-Bent (1) \\
Bright-core & 3 & 3 & -- \\
FRI & 3 & 2 & Hybrid (1) \\
Core-Displaced & 3 & 3 & -- \\
FRII & 3 & 3 & -- \\
Hybrid & 3 & 3 & -- \\
Large FRII & 3 & 3 & -- \\
Spirals & 3 & 3 & -- \\
Structured-Bent & 3 & 1 & ACJ (2) \\
Unresolved & 3 & 3 & -- \\
\bottomrule
\end{tabular}

\caption[Test-set classification counts]{Predictions for the \(33\) test sources.}
\label{tab:heldout_counts}
\end{table}

The \(F_1\) uncertainties in the main comparison tables and Table~\ref{tab:las_sensitivity} use source-level bootstrap resampling: for each catalogue separately, we repeatedly draw samples of the original size, allowing a source to be selected more than once. Each selected source retains its reference label and saved model prediction, and we recalculate \(F_1\) for each sample without retraining the model. The central \(95\,\%\) of the resulting \(F_1\) values defines the reported uncertainty interval.

Signed uncertainties give the distances from each estimate to its lower and upper \(95\,\%\) interval bounds. They describe sampling variation within the selected samples, excluding uncertainty in reference labels, source selection and training. FRII includes both internal FRII classes, as in the main comparisons.

\begin{table}[H]
\centering
\small
\setstretch{1}
\setlength{\tabcolsep}{5pt}
\renewcommand{\arraystretch}{1.6}
\begin{tabular}{@{}llrrr@{}}
\toprule
Catalogue & LAS (arcsec) & Sources & FRI \(F_1\) & FRII \(F_1\) \\
\midrule
Clews & 60--120 & 6002 & \(0.310_{-0.024}^{+0.024}\) & \(0.739_{-0.014}^{+0.013}\) \\
Horton & 60--120 & 4635 & \(0.120_{-0.025}^{+0.026}\) & \(0.772_{-0.012}^{+0.011}\) \\
Chilufya & 60--120 & 1420 & \(0.388_{-0.059}^{+0.057}\) & \(0.378_{-0.046}^{+0.044}\) \\
Clews & 120--240 & 1539 & \(0.562_{-0.036}^{+0.035}\) & \(0.739_{-0.032}^{+0.029}\) \\
Horton & 120--240 & 2868 & \(0.383_{-0.032}^{+0.033}\) & \(0.719_{-0.019}^{+0.018}\) \\
Chilufya & 120--240 & 906 & \(0.490_{-0.051}^{+0.050}\) & \(0.521_{-0.055}^{+0.051}\) \\
Clews & 240--1000 & 216 & \(0.419_{-0.105}^{+0.093}\) & \(0.583_{-0.099}^{+0.092}\) \\
Horton & 240--1000 & 849 & \(0.363_{-0.055}^{+0.054}\) & \(0.605_{-0.045}^{+0.042}\) \\
Chilufya & 240--1000 & 264 & \(0.439_{-0.101}^{+0.093}\) & \(0.592_{-0.089}^{+0.081}\) \\
\bottomrule
\end{tabular}

\caption[Angular-size sensitivity of external correspondence]{Curated-excluded comparisons by LAS, with \(95\,\%\) bootstrap intervals. Bins include their lower edge; the final bin includes \(1000''\).}
\label{tab:las_sensitivity}
\end{table}

Removing Clews's uncertain FR labels as well as all curated sources preserves the stronger FRII correspondence. This checks sensitivity to reference selection without resolving the remaining mapping ambiguities.
\par
\endgroup
